\section{Newsvendor Problem and Multi-Echelon Extensions}

The newsvendor problem~\citep{arrow1951} determines the optimal order quantity $Q^* = F^{-1}\!\left(\frac{p-c}{p-s}\right)$ for a single perishable product under stochastic demand~\citep{porteus2002}. Clark and Scarf~\citep{clark1960} extended this to multi-echelon systems, with later generalisations by Federgruen and Zipkin~\citep{federgruen1984} and Graves and Willems~\citep{graves2000}. However, analytical tractability diminishes with scale, motivating simulation-based approaches.

\section{Monte-Carlo Simulation and Correlated Demand}

Monte-Carlo methods with sample average approximation (SAA) offer flexibility for complex, correlated demand~\citep{birge2011, shapiro2014}. Spatial correlations are modelled via the Cholesky decomposition $\bm{\Sigma} = \mathbf{L}\mathbf{L}^\top$, generating correlated scenarios as $\mathbf{D} = \bm{\mu} + \mathbf{L}\mathbf{Z}$~\citep{glasserman2003}. The M5 forecasting dataset~\citep{makridakis2022} provides realistic hierarchical correlation structures across 3,049 products and 10 stores.

\section{GPU Computing and the Memory Wall}

Modern GPUs offer massive parallelism but suffer from the ``memory wall''~\citep{wulf1995}. Ivanov et al.~\citep{ivanov2021} showed that intermediate tensor materialisation dominates cost in data-intensive workloads. For the newsvendor problem, the $\mathbf{D}[N,S]$ matrix must make a round-trip through HBM, creating the primary bottleneck.

\section{Kernel Fusion, Triton, and FlashAttention}

Kernel fusion eliminates intermediate memory accesses by combining sequential GPU kernels~\citep{wang2010}. The Triton language~\citep{tillet2019} provides a Python DSL for writing custom tiled GPU kernels with autotuning. FlashAttention~\citep{dao2022, dao2023} demonstrated that fusing matmul with downstream operations within SRAM avoids materialising the $N \times N$ attention matrix---a direct analogue to our approach. PyTorch's \texttt{torch.compile}~\citep{ansel2024} can fuse element-wise operations but \emph{cannot} fuse across the cuBLAS matmul boundary, leaving the $\mathbf{D}$ matrix in HBM.

\section{Risk Measures and Related Work}

Conditional Value-at-Risk (CVaR), also known as Expected Shortfall, was formalised by Rockafellar and Uryasev~\citep{rockafellar2000} as a coherent risk measure that captures tail behaviour beyond Value-at-Risk (VaR). For a confidence level $\alpha \in (0,1)$, $\text{CVaR}_\alpha$ is defined as the expected value of the loss distribution in the worst $(1-\alpha)$ fraction of scenarios. Formally, if $\pi$ denotes profit, $\text{CVaR}_\alpha(\pi) = \mathbb{E}[\pi \mid \pi \leq q_\alpha]$, where $q_\alpha = F^{-1}(1-\alpha)$ is the $\alpha$-quantile of the profit distribution. Computing CVaR exactly requires sorting all $S$ scenario profits to identify the quantile threshold---an $O(S \log S)$ operation that is particularly expensive on GPUs, where comparison-based sorting requires multiple synchronisation barriers across thread blocks. Histogram-based approximations offer a practical alternative: by binning the profit values into $B$ equally spaced buckets and accumulating counts, the quantile can be located by scanning the cumulative histogram, reducing the problem to $O(S + B)$ with embarrassingly parallel binning. This histogram approach is well-suited to the GPU's SIMT execution model and avoids the global synchronisation overhead of sorting.

CVaR has been applied to newsvendor settings by Gotoh and Takano~\citep{gotoh2007} and Chen et al.~\citep{chen2009}, who demonstrated that risk-averse ordering policies reduce the probability of catastrophic losses at the expense of modest reductions in expected profit. Prior GPU work on inventory models focused on parallelising scenario evaluation across CUDA cores~\citep{pichitlamken2012} or accelerating the underlying linear algebra with cuBLAS~\citep{munguia2019}, but these approaches still materialise intermediate tensors in HBM. To our knowledge, fusing the demand matmul with newsvendor logic into a single SRAM-resident kernel has not been previously reported in the literature. The approach is conceptually analogous to FlashAttention's treatment of the attention score matrix as a transient intermediate, but applied to the inventory optimisation domain where the intermediate demand matrix $\mathbf{D}$ plays the analogous role.

\section{Summary of Research Gaps}

The literature reveals three key gaps that this thesis addresses.

\textbf{Gap 1: Kernel fusion for inventory optimisation.} Existing GPU-accelerated inventory models parallelise scenario evaluation across CUDA cores but still materialise intermediate tensors---most critically the demand matrix $\mathbf{D} \in \mathbb{R}^{N \times S}$---in high-bandwidth memory (HBM). This materialisation creates a round-trip memory access pattern where the demand matrix is written to HBM after the matmul and then read back for profit computation, consuming bandwidth that dominates total execution time. The kernel fusion paradigm pioneered by FlashAttention~\citep{dao2022} for transformer attention has not been applied to the inventory optimisation domain, despite the structural similarity between the transient attention-score matrix and the transient demand matrix.

\textbf{Gap 2: The \texttt{torch.compile} fusion boundary.} PyTorch's \texttt{torch.compile} with the Inductor backend~\citep{ansel2024} can fuse chains of element-wise operations into a single kernel, but it cannot fuse operations across the cuBLAS matrix multiplication boundary. The demand generation step $\mathbf{D} = \bm{\mu} + \mathbf{L}\mathbf{Z}$ invokes cuBLAS for the $\mathbf{L}\mathbf{Z}$ matmul, and the subsequent newsvendor profit logic (clamping, min/max, multiply-accumulate) is compiled into a separate kernel. This two-kernel pattern forces the full $\mathbf{D}$ matrix through HBM, negating much of the benefit of GPU parallelism for memory-bound problem sizes. Closing this gap requires a custom kernel that implements both the tiled matmul and the downstream profit arithmetic within a single SRAM-resident pass.

\textbf{Gap 3: Risk-aware and constrained GPU-fused extensions.} While CVaR-based newsvendor formulations exist in the operations research literature~\citep{gotoh2007, chen2009} and budget-constrained inventory models have been studied analytically~\citep{porteus2002}, GPU-fused implementations of these extensions have not been explored. Computing CVaR on GPU requires efficient tail-quantile estimation over millions of scenarios, and budget constraints introduce coupling across nodes that complicates parallelisation. Developing these extensions within the fused-kernel framework demonstrates both the generality of the approach and its practical relevance for risk-aware supply chain decision-making.

The Triton programming model provides the ideal platform to close these gaps, offering Python-level accessibility with hardware-level control over SRAM tiling and autotuning across GPU architectures.

The ideal solution would combine the tiled matmul capability of cuBLAS with the element-wise fusion of the Inductor backend, all within a single kernel that respects the SRAM budget of the target GPU. This is precisely what the Triton programming model enables: the developer specifies the tile dimensions and the fused computation in Python, and the Triton compiler generates machine code that manages shared memory, register allocation, and memory coalescing automatically. The resulting kernel can be autotuned across a small grid of tile configurations, yielding near-optimal performance without manual assembly-level optimisation.
